% =====================================================================
% Kapitel 8: OpenTUI-Explorer
% =====================================================================

\chapter{OpenTUI-Explorer}

\label{ch:tui}

Der \emph{OpenTUI-Explorer} ist die schnellste und interaktivste
Möglichkeit, die gesammelten Daten zu durchsuchen. Er läuft direkt
im Terminal, ohne Jupyter, ohne Notebook, ohne Python-Overhead im
Hot-Path.

\section{Was ist OpenTUI?}

\begin{definition}
\textbf{OpenTUI}
OpenTUI (https://opentui.com) ist ein von Anomaly (https://anoma.ly)
entwickeltes TUI-Framework. Es besteht aus:

\begin{itemize}
  \item einem in \emph{Zig} geschriebenen \emph{Core}
        (schnell, plattformübergreifend, nativ),
  \item einer \emph{C-ABI} für Bindings aus jeder Sprache,
  \item \emph{TypeScript}-Bindings mit React- und Solid.js-Support.
\end{itemize}

Der Kern nutzt \emph{Yoga} (Facebooks Flexbox-Engine) für Layouts
und rendert das Ergebnis direkt in einen Terminal-Buffer.
\end{definition}

\begin{merke}
OpenTUI ist \emph{die gleiche Engine}, die \foreign{opencode}
(powered by Anomaly) und \foreign{terminal.shop} antreibt. Es
befindet sich in aktiver Entwicklung, kann also Breaking Changes
enthalten.
\end{merke}

\section{Warum TUI statt GUI?}

\begin{figure}[ht]
\centering
\begin{tabular}{lll}
  \toprule
  \textbf{Aspekt} & \textbf{TUI} & \textbf{GUI (z.\,B.\ Web-App)} \\
  \midrule
  Startup-Zeit          & $<$100 ms (Bun)         & $>$1 s (Browser) \\
  Speicherverbrauch     & $\sim$50 MB               & $\sim$300 MB \\
  Tastatur-Effizienz    & \foreign{Vim-style}      & Maus \\
  SSH-/Server-tauglich  & ja                       & nur mit X-Forwarding \\
  Visueller Reiz        & niedrig                  & hoch \\
  Lernkurve             & hoch (Tasten lernen)     & niedrig (klicken) \\
  \bottomrule
\end{tabular}
\caption{TUI vs.\ GUI}
\label{fig:tui-vs-gui}
\end{figure}

Für die tägliche Exploration von Finanzmarktdaten -- springen
zwischen 12~Serien, monatliche Updates prüfen, Daten vergleichen --
ist eine TUI das richtige Werkzeug.

\section{Architektur des TUI-Layers}

Das TUI ist ein \emph{eigenständiger TypeScript-Subtree} im
Projekt-Root:

\begin{lstlisting}[language=shell, caption={TUI-Verzeichnisstruktur}]
tui/
├── package.json            # Bun + React 19 + OpenTUI 0.4
├── tsconfig.json
├── bun.lock
├── src/
│   ├── index.tsx          # Entry: createCliRenderer + createRoot
│   ├── App.tsx            # Root-Komponente (Layout, State, Keyboard)
│   ├── components/
│   │   ├── SeriesList.tsx   # Linke Seite: 12 Serien
│   │   └── SeriesDetail.tsx # Rechte Seite: Stats + Tail
│   ├── data/
│   │   ├── db.ts          # pg-Pool
│   │   └── queries.ts     # 4 SQL-Queries
│   └── types.ts           # TypeScript-Typen
├── tests/                 # bun test (Unit + Integration)
└── scripts/check-db.ts    # DB-Connectivity-Test
\end{lstlisting}

\section{Die Komponenten}

\subsection{Entry Point: \code{index.tsx}}

\begin{lstlisting}[language=typescript, caption={TUI-Entry-Point},label=lst:tui-entry]
import { createCliRenderer } from "@opentui/core";
import { createRoot } from "@opentui/react";
import { App } from "./App";
import { closePool } from "./data/db";

const renderer = await createCliRenderer({ exitOnCtrlC: true });
const root = createRoot(renderer);
root.render(<App />);
\end{lstlisting}

Der Entry-Point ist minimal:

\begin{enumerate}
  \item \code{createCliRenderer()} -- startet den OpenTUI-Renderer,
        der den Terminal-Buffer verwaltet.
  \item \code{createRoot(renderer)} -- erzeugt einen React-Root,
        der direkt in den OpenTUI-Buffer rendert.
  \item \code{root.render(<App />)} -- mountet die App-Komponente.
\end{enumerate}

\subsection{Root: \code{App.tsx}}

Die Root-Komponente verwaltet drei Dinge: \emph{State} (welche Serien
sind in der DB, welche ist selektiert), \emph{Keyboard-Handling}
und das \emph{Layout} (2-Spalten).

\begin{lstlisting}[language=typescript, caption={TUI-App: State + Keyboard},label=lst:tui-app]
function App() {
  const renderer = useRenderer();
  const [series, setSeries] = useState<SeriesRow[]>([]);
  const [selectedIndex, setSelectedIndex] = useState(0);

  useEffect(() => {
    listSeries().then(setSeries).catch(/* ... */);
  }, []);

  useKeyboard((key) => {
    if (key.name === "q") renderer.destroy();
    if (key.name === "r") setRefreshToken(t => t + 1);
    if (key.name === "up" || key.name === "k") setSelectedIndex(i => Math.max(0, i - 1));
    if (key.name === "down" || key.name === "j") setSelectedIndex(i => Math.min(series.length - 1, i + 1));
  });

  return (
    <box flexDirection="row">
      <SeriesList series={series} selectedIndex={selectedIndex} />
      {series[selectedIndex] && <SeriesDetail series={series[selectedIndex]} />}
    </box>
  );
}
\end{lstlisting}

\subsection{SeriesList -- Linke Spalte}

Die \code{SeriesList}-Komponente rendert alle geladenen Serien als
\emph{scrollbare} Liste. Jede Zeile enthält:

\begin{itemize}
  \item \textbf{Badge}: \code{M} für Macro (FRED/ECB), \code{E} für
        Equity (Yahoo).
  \item \textbf{Series-ID}: z.\,B.\ \code{FRED:DGS3MO} oder
        \code{ECB:EXR:USD.EUR.SP00.A}.
  \item \textbf{Vendor-Tag}: \code{fred}, \code{ecb}, \code{yahoo}.
  \item \textbf{Row-Count}: rechtsbündig für schnelles Scannen.
\end{itemize}

\subsection{SeriesDetail -- Rechte Spalte}

Die \code{SeriesDetail}-Komponente zeigt für die ausgewählte Serie:

\begin{itemize}
  \item \textbf{Header}: Series-ID + Name.
  \item \textbf{Stats-Panel}: rows, first/last date, unit, source,
        min/max/mean.
  \item \textbf{Tail-Tabelle}: die letzten 30~Werte, bei Macro
        als \code{date, value}, bei Equity als
        \code{date, close, adj\_close, vol}.
\end{itemize}

\section{Die SQL-Queries}

Das TUI spricht \emph{direkt} mit der PostgreSQL-DB via
\code{node-postgres} -- es gibt keinen Python-Overhead im Hot-Path.

\begin{lstlisting}[language=typescript, caption={4 SQL-Queries im TUI}]
const LIST_SERIES_SQL = `
  SELECT 'macro' AS kind, series_id AS id, name, source, frequency, unit, COUNT(*)::int AS n
  FROM macro_observations mo
  JOIN macro_series ms USING (series_id)
  GROUP BY series_id, name, source, frequency, unit
  UNION ALL
  SELECT 'equity' AS kind, COALESCE(ii.value, i.instrument_id::text) AS id,
         i.name, 'yahoo' AS source, 'D' AS frequency, i.currency AS unit, COUNT(*)::int AS n
  FROM prices_daily pd
  JOIN instruments i USING (instrument_id)
  LEFT JOIN instrument_identifiers ii
    ON ii.instrument_id = i.instrument_id AND ii.scheme = 'YAHOO'
  GROUP BY i.instrument_id, ii.value, i.name, i.currency
  ORDER BY 1, 2
`;

const MACRO_TAIL_SQL = `
  SELECT to_char(date, 'YYYY-MM-DD') AS date, value
  FROM macro_observations WHERE series_id = $1
  ORDER BY date DESC LIMIT $2
`;

const EQUITY_TAIL_SQL = `
  SELECT to_char(pd.date, 'YYYY-MM-DD') AS date, pd.close, pd.adjusted_close, pd.volume
  FROM prices_daily pd
  JOIN instrument_identifiers ii
    ON ii.instrument_id = pd.instrument_id AND ii.scheme = 'YAHOO'
  WHERE ii.value = $1
  ORDER BY pd.date DESC LIMIT $2
`;

const MACRO_STATS_SQL = `
  SELECT COUNT(*)::int AS n, to_char(MIN(date), 'YYYY-MM-DD') AS first,
         to_char(MAX(date), 'YYYY-MM-DD') AS last, MIN(value) AS min, MAX(value) AS max, AVG(value) AS mean
  FROM macro_observations WHERE series_id = $1
`;
\end{lstlisting}

Beachten Sie:

\begin{itemize}
  \item \textbf{UNION ALL} -- eine einzige Query listet Macro- und
        Equity-Serien in einer Liste.
  \item \code{COALESCE(ii.value, i.instrument\_id::text)} -- für
        Equity-Serien wird der YAHOO-Ticker (\code{AAPL}) statt
        der numerischen \code{instrument\_id} (\code{1}) angezeigt.
  \item \code{LEFT JOIN instrument\_identifiers} -- damit auch
        Equity-Serien ohne YAHOO-Identifier angezeigt würden (was
        aktuell nicht vorkommt, aber defensiv ist).
\end{itemize}

\section{Tastatur-Handling}

\begin{table}[ht]
\centering
\small
\begin{tabular}{ll}
  \toprule
  \textbf{Taste} & \textbf{Aktion} \\
  \midrule
  \code{$\uparrow$} / \code{k}  & Selection nach oben \\
  \code{$\downarrow$} / \code{j} & Selection nach unten \\
  \code{g} / \code{Home}  & Zum Anfang \\
  \code{G} / \code{End}   & Zum Ende \\
  \code{r}                & Re-Query (für Live-Refresh) \\
  \code{q} / \code{Ctrl-C}  & Quit \\
  \bottomrule
\end{tabular}
\caption{Tastatur-Bindings des TUI. Die \code{j}/\code{k}-Variante
  ist Vim-style für Entwickler, die diese Tasten gewohnt sind.}
\label{tab:tui-keys}
\end{table}

\begin{lstlisting}[language=typescript, caption={Keyboard-Handler},label=lst:tui-keys]
useKeyboard((key) => {
  if (key.name === "q" || (key.ctrl && key.name === "c")) {
    renderer.destroy();
    setTimeout(() => process.exit(0), 50);
    return;
  }
  if (key.name === "r") {
    setRefreshToken(t => t + 1);
    return;
  }
  if (key.name === "up" || key.name === "k") {
    setSelectedIndex(i => Math.max(0, i - 1));
    return;
  }
  // ... etc.
});
\end{lstlisting}

\section{Tests}

Das TUI hat 27~Tests in 4~Dateien:

\begin{itemize}
  \item \code{src/data/queries.test.ts} (9~Tests) -- SQL-Strings,
        Row-Mapping, Param-Übergabe.
  \item \code{src/data/db.test.ts} (3~Tests) -- Public-API des
        \code{db}-Moduls.
  \item \code{src/types/format.test.ts} (8~Tests) -- Pure-Funktionen
        (\code{fmtNumber}, \code{fmtPct}).
  \item \code{tests/integration.test.ts} (7~Tests) -- Echtes
        Postgres, 12~Series, OHLCV-Validierung.
\end{itemize}

\begin{lstlisting}[language=bash, caption={TUI-Tests laufen}]
cd tui && bun test
# 27 pass, 0 fail, 203 expect(), 4 files
\end{lstlisting}

\section{Was als Nächstes?}

Die TUI ist die \emph{operative Oberfläche} des Projekts. Das
nächste Kapitel zeigt, wie wir die Datenqualität \emph{testen} --
nicht nur funktional (geht der Code?), sondern auch inhaltlich
(stimmen die Daten?).
